\section{Limitations and Future Work}
\label{sec:limitations}

Several limitations qualify the interpretation of our results.
First, the model was trained exclusively on South Korean NHIS data; differences in diet, genetic predisposition, and healthcare access across populations mean that discrimination and calibration may not transfer to other countries without retraining~\cite{Sun2022idf}.
Second, the diabetes label is a proxy derived from a single fasting plasma glucose reading (FPG $\geq$ 126~mg/dL) rather than a confirmed clinical diagnosis.
The American Diabetes Association also recognizes HbA1c $\geq$ 6.5\% as a diagnostic criterion~\cite{ADA2024standards}, but HbA1c is not collected in the NHIS general health examination.
This proxy cannot distinguish newly undiagnosed cases from patients with known but poorly controlled T2D, introducing label noise in both directions.
Third, the cross-sectional design identifies individuals who likely have diabetes at the time of the examination; it does not predict future diabetes onset, limiting its use for incident-risk stratification.
Fourth, operating at the recall-oriented threshold produces a 37.5\% false-positive rate (1 $-$ specificity).
Although each false positive triggers only a confirmatory blood test, the downstream cost of additional testing may be substantial when deployed at national scale.
Additionally, all reported metrics derive from a single random seed (42); we did not perform repeated splits or statistical significance testing, so the observed AUC differences between models may fall within sampling variability.
We also did not conduct subgroup analyses stratified by age, sex, or BMI category, which would be necessary to assess whether model performance is equitable across clinically relevant subpopulations.
Finally, patients who lack optional laboratory values may differ systematically from those with complete records---the no-lipid subgroup tends to be younger and healthier---so predictions made with minimal inputs may underestimate risk for atypical patients within that subgroup.

Several directions could address these limitations.
Multi-national validation using cohorts from countries with different demographic and dietary profiles would clarify the model's transferability~\cite{Lee2025multicohort}.
Incorporating longitudinal NHIS records would enable prediction of incident diabetes rather than detection of existing cases.
Integrating the model's calibrated risk score with established screening guidelines such as the USPSTF recommendation~\cite{USPSTF2021screening} or the FINDRISC questionnaire could produce hybrid protocols that leverage both clinical heuristics and data-driven pattern recognition.
On the modeling side, the nine binary missingness indicator flags contributed negligible SHAP importance and could be removed to simplify the feature set without expected performance loss.
Regional calibration of probability thresholds for populations with different baseline prevalence rates would also improve the clinical utility of the risk tiers when the model is applied outside South Korea~\cite{Sun2022idf}.
